\documentclass[../../main.tex]{subfiles}
\graphicspath{{../../image/}{../../figure/}}
\begin{document}

\section{松弛法}\label{section:ch06-松弛法}

构造\eqref{align:ch06-6-1-1}可容许弱解的另一个有趣方法是通过\emph{松弛}。出发点是一个含 $m+1$ 个方程的半线性组：
\begin{align}
\begin{cases}
\partial_{t}v(x,t)+\displaystyle\sum_{\alpha=1}^{m}c_{\alpha}\partial_{\alpha}v(x,t)=\dfrac{1}{\mu}\displaystyle\sum_{\alpha=1}^{m}[F_{\alpha}(v(x,t))-Z_{\alpha}(x,t)],\\[8pt]
\partial_{t}Z_{\alpha}(x,t)-c_{\alpha}\partial_{\alpha}Z_{\alpha}(x,t)=\dfrac{1}{\mu}[F_{\alpha}(v(x,t))-Z_{\alpha}(x,t)],\qquad \alpha=1,\cdots,m,
\end{cases}
\label{align:ch06-6-6-1}
\end{align}
未知量为 $(v,Z_{1},\cdots,Z_{m})$，$\mu$ 为小正参数，$c_{\alpha}$ 为给定常数，$F_{\alpha}$ 为指定光滑函数，满足
\begin{align}
F_{\alpha}^{\prime}(v)<0,\qquad -\infty<v<\infty,\ \alpha=1,\cdots,m,
\label{align:ch06-6-6-2}
\end{align}
\begin{align}
F_{\alpha}(0)=0,\qquad F_{\alpha}(v)\to\pm\infty,\ \text{当 }v\to\mp\infty,\ \alpha=1,\cdots,m.
\label{align:ch06-6-6-3}
\end{align}

注意\eqref{align:ch06-6-6-1}的解满足守恒律
\begin{align}
\partial_{t}\Bigl{[}v(x,t)-\sum_{\alpha=1}^{m}Z_{\alpha}(x,t)\Bigr{]}+\sum_{\alpha=1}^{m}c_{\alpha}\partial_{\alpha}\Bigl{[}v(x,t)+Z_{\alpha}(x,t)\Bigr{]}=0.
\label{align:ch06-6-6-4}
\end{align}
由\eqref{align:ch06-6-6-1}右端的形式可预期，当 $\mu\downarrow 0$ 时变量 $Z_{\alpha}$“松弛”到其平衡态 $F_{\alpha}(v)$；此时\eqref{align:ch06-6-6-4}化为标量守恒律\eqref{align:ch06-6-1-1}，其中
\begin{align}
u=v-\sum_{\alpha=1}^{m}F_{\alpha}(v),\qquad G_{\alpha}(u)=c_{\alpha}[v+F_{\alpha}(v)],\ \alpha=1,\cdots,m.
\label{align:ch06-6-6-5}
\end{align}

上述考虑提示一个方案：把\eqref{align:ch06-6-1-1}的解构造为\eqref{align:ch06-6-6-1}解的渐近极限。

第一步是考察\eqref{align:ch06-6-6-1}在指定初值下的 Cauchy 问题：
\begin{align}
v(x,0)=v_{0}(x),\qquad Z_{\alpha}(x,0)=Z_{\alpha 0}(x),\ \alpha=1,\cdots,m,\ x\in\mathbb{R}^{m}.
\label{align:ch06-6-6-6}
\end{align}
由于\eqref{align:ch06-6-6-1}是半线性双曲组，当初值 $(v_{0},Z_{10},\cdots,Z_{m0})\in C_{0}^{1}(\mathbb{R}^{m})$ 时，在最大时间区间 $[0,T)$（$0<T\leqslant\infty$）上存在唯一经典解 $(v,Z_{1},\cdots,Z_{m})$。对任一 $t\in[0,T)$，函数 $(v(\cdot,t),Z_{1}(\cdot,t),\cdots,Z_{m}(\cdot,t))\in C_{0}^{1}(\mathbb{R}^{m})$。此外，若 $T<\infty$，则
\begin{align}
\|v(\cdot,t)\|_{L^{\infty}(\mathbb{R}^{m})}+\sum_{\alpha=1}^{m}\|Z_{\alpha}(\cdot,t)\|_{L^{\infty}(\mathbb{R}^{m})}\to\infty,\qquad\text{当 }t\uparrow T.
\label{align:ch06-6-6-7}
\end{align}

这里需要更广泛初值类下的（可能是弱的）解，且在时间上整体存在。在假设\eqref{align:ch06-6-6-2}、\eqref{align:ch06-6-6-3}下，这样的解确实存在，因为\eqref{align:ch06-6-6-1}右端的效应是耗散的。这体现于下述定理。

\begin{theorem}\label{theorem:ch06-6-6-1}
对任何初值 $(v_{0},Z_{10},\cdots,Z_{m0})\in L^{1}(\mathbb{R}^{m})\cap L^{\infty}(\mathbb{R}^{m})$，\eqref{align:ch06-6-6-1}、\eqref{align:ch06-6-6-6}在 $\mathbb{R}^{m}\times[0,\infty)$ 上存在唯一弱解 $(v,Z_{1},\cdots,Z_{m})$，使 $(v(\cdot,t),Z_{1}(\cdot,t),\cdots,Z_{m}(\cdot,t))\in C^{0}([0,\infty);L^{1}(\mathbb{R}^{m}))$。若
\begin{align}
a\leqslant v_{0}(x)\leqslant b,\qquad F_{\alpha}(b)\leqslant Z_{\alpha 0}(x)\leqslant F_{\alpha}(a),\ \alpha=1,\cdots,m,\ x\in\mathbb{R}^{m},
\label{align:ch06-6-6-8}
\end{align}
则
\begin{align}
a\leqslant v(x,t)\leqslant b,\qquad F_{\alpha}(b)\leqslant Z_{\alpha}(x,t)\leqslant F_{\alpha}(a),\ \alpha=1,\cdots,m,\ (x,t)\in\mathbb{R}^{m}\times[0,\infty).
\label{align:ch06-6-6-9}
\end{align}
此外，若 $(\bar{v},\bar{Z}_{1},\cdots,\bar{Z}_{m})$ 是\eqref{align:ch06-6-6-1}的另一解，初值 $(\bar{v}_{0},\bar{Z}_{10},\cdots,\bar{Z}_{m0})\in L^{1}(\mathbb{R}^{m})\cap L^{\infty}(\mathbb{R}^{m})$，则对任意 $t\in[0,\infty)$，
\begin{align}
&\int_{\mathbb{R}^{m}}\Bigl{\{}[v(x,t)-\bar{v}(x,t)]^{+}+\sum_{\alpha=1}^{m}[\bar{Z}_{\alpha}(x,t)-Z_{\alpha}(x,t)]^{+}\Bigr{\}}\,\mathrm{d}x\notag\\
&\qquad\leqslant\int_{\mathbb{R}^{m}}\Bigl{\{}[v_{0}(x)-\bar{v}_{0}(x)]^{+}+\sum_{\alpha=1}^{m}[\bar{Z}_{\alpha 0}(x)-Z_{\alpha 0}(x)]^{+}\Bigr{\}}\,\mathrm{d}x,
\label{align:ch06-6-6-10}
\end{align}
\begin{align}
&\|v(\cdot,t)-\bar{v}(\cdot,t)\|_{L^{1}(\mathbb{R}^{m})}+\sum_{\alpha=1}^{m}\|Z_{\alpha}(\cdot,t)-\bar{Z}_{\alpha}(\cdot,t)\|_{L^{1}(\mathbb{R}^{m})}\notag\\
&\qquad\leqslant\|v_{0}(\cdot)-\bar{v}_{0}(\cdot)\|_{L^{1}(\mathbb{R}^{m})}+\sum_{\alpha=1}^{m}\|Z_{\alpha 0}(\cdot)-\bar{Z}_{\alpha 0}(\cdot)\|_{L^{1}(\mathbb{R}^{m})}.
\label{align:ch06-6-6-11}
\end{align}
特别地，若
\begin{align}
v_{0}(x)\leqslant\bar{v}_{0}(x),\qquad Z_{\alpha 0}(x)\geqslant\bar{Z}_{\alpha 0}(x),\ \alpha=1,\cdots,m,\ x\in\mathbb{R}^{m},
\label{align:ch06-6-6-12}
\end{align}
则
\begin{align}
v(x,t)\leqslant\bar{v}(x,t),\qquad Z_{\alpha}(x,t)\geqslant\bar{Z}_{\alpha}(x,t),\ \alpha=1,\cdots,m,\ (x,t)\in\mathbb{R}^{m}\times[0,\infty).
\label{align:ch06-6-6-13}
\end{align}
\end{theorem}

\begin{proof}
第一步在简化假设下建立\eqref{align:ch06-6-6-10}：两解 $(v,Z_{1},\cdots,Z_{m})$ 与 $(\bar{v},\bar{Z}_{1},\cdots,\bar{Z}_{m})$ 均为经典解，初值 $(v_{0},Z_{10},\cdots,Z_{m0})$、$(\bar{v}_{0},\bar{Z}_{10},\cdots,\bar{Z}_{m0})\in C_{0}^{1}(\mathbb{R}^{m})$。对 $\varepsilon>0$，回顾由\eqref{align:ch06-6-3-8}定义的函数 $\eta_{\varepsilon}$，注意到
\begin{align}
&\partial_{t}\Bigl{[}\eta_{\varepsilon}(v-\bar{v})+\sum_{\alpha=1}^{m}\eta_{\varepsilon}(\bar{Z}_{\alpha}-Z_{\alpha})\Bigr{]}+\sum_{\alpha=1}^{m}c_{\alpha}\partial_{\alpha}\Bigl{[}\eta_{\varepsilon}(v-\bar{v})-\eta_{\varepsilon}(\bar{Z}_{\alpha}-Z_{\alpha})\Bigr{]}\notag\\
&\qquad=\frac{1}{\mu}\sum_{\alpha=1}^{m}[\eta_{\varepsilon}^{\prime}(v-\bar{v})-\eta_{\varepsilon}^{\prime}(\bar{Z}_{\alpha}-Z_{\alpha})][F_{\alpha}(v)-F_{\alpha}(\bar{v})+\bar{Z}_{\alpha}-Z_{\alpha}]
\label{align:ch06-6-6-14}
\end{align}
直接由\eqref{align:ch06-6-6-1}得到。对 $v,\bar{v},Z_{\alpha},\bar{Z}_{\alpha}$ 的固定值（无论符号），\eqref{align:ch06-6-6-14}右端当 $\varepsilon\downarrow 0$ 时具有非正极限。于是在 $\mathbb{R}^{m}\times(0,t)$ 上积分\eqref{align:ch06-6-6-14}并令 $\varepsilon\downarrow 0$，即得\eqref{align:ch06-6-6-10}。

当\eqref{align:ch06-6-6-12}成立时，\eqref{align:ch06-6-6-10}立即蕴含\eqref{align:ch06-6-6-13}。注意对任意常数 $a$、$b$，$(a,F_{1}(a),\cdots,F_{m}(a))$ 与 $(b,F_{1}(b),\cdots,F_{m}(b))$ 都是\eqref{align:ch06-6-6-1}的特解，因而\eqref{align:ch06-6-6-8}蕴含\eqref{align:ch06-6-6-9}。特别地，不可能发生 blow-up\eqref{align:ch06-6-6-7}，故解在 $\mathbb{R}^{m}\times[0,\infty)$ 上整体存在。

为得\eqref{align:ch06-6-6-11}，只需在\eqref{align:ch06-6-6-10}中交换 $(v,Z_{1},\cdots,Z_{m})$ 与 $(\bar{v},\bar{Z}_{1},\cdots,\bar{Z}_{m})$ 的地位，再把所得不等式与原\eqref{align:ch06-6-6-10}相加。

至此我们在初值属于 $C_{0}^{1}(\mathbb{R}^{m})$ 的经典解框架下验证了定理的全部论断。然而，凭借 L$^{1}$ 压缩估计\eqref{align:ch06-6-6-11}，以 $L^{1}(\mathbb{R}^{m})\cap L^{\infty}(\mathbb{R}^{m})$ 中任意初值满足所述性质的\eqref{align:ch06-6-6-1}弱解，可轻易由经典解序列的 L$^{1}$ 极限构造。证明完成。
\end{proof}

下一步考察\eqref{align:ch06-6-6-1}的解当 $\mu\downarrow 0$ 时的极限行为。诱导 $Z_{\alpha}$ 松弛到平衡值 $F_{\alpha}(v)$ 的机制将通过一个类熵不等式来刻画。我们定义族
\begin{align}
\Phi_{\alpha}(Z_{\alpha})=-\int_{0}^{Z_{\alpha}}F_{\alpha}^{-1}(w)\,\mathrm{d}w,\qquad \alpha=1,\cdots,m,
\label{align:ch06-6-6-15}
\end{align}
它是 $(-\infty,\infty)$ 上非负凸函数。设 $(v,Z_{1},\cdots,Z_{m})$ 是\eqref{align:ch06-6-6-1}的经典解，初值 $(v_{0},Z_{10},\cdots,Z_{m0})\in C_{0}^{1}(\mathbb{R}^{m})$，则易验证
\begin{align}
\partial_{t}\Bigl{[}\tfrac{1}{2}v^{2}+\sum_{\alpha=1}^{m}\Phi_{\alpha}(Z_{\alpha})\Bigr{]}+\sum_{\alpha=1}^{m}c_{\alpha}\partial_{\alpha}\Bigl{[}\tfrac{1}{2}v^{2}-\Phi_{\alpha}(Z_{\alpha})\Bigr{]}
=\frac{1}{\mu}\sum_{\alpha=1}^{m}[v-F_{\alpha}^{-1}(Z_{\alpha})][F_{\alpha}(v)-Z_{\alpha}].
\label{align:ch06-6-6-16}
\end{align}
由于 $v-F_{\alpha}^{-1}(Z_{\alpha})=F_{\alpha}^{-1}(F_{\alpha}(v))-F_{\alpha}^{-1}(Z_{\alpha})$，中值定理蕴含
\begin{align}
-[v-F_{\alpha}^{-1}(Z_{\alpha})][F_{\alpha}(v)-Z_{\alpha}]\geqslant\frac{1}{k}[F_{\alpha}(v)-Z_{\alpha}]^{2},
\label{align:ch06-6-6-17}
\end{align}
其中 $k$ 是 $-F_{\alpha}^{\prime}$ 在 $v$ 的值域上的任一上界。于是在 $\mathbb{R}^{m}\times[0,\infty)$ 上积分\eqref{align:ch06-6-6-16}，得不等式
\begin{align}
\int_{0}^{\infty}\int_{\mathbb{R}^{m}}\sum_{\alpha=1}^{m}[F_{\alpha}(v)-Z_{\alpha}]^{2}\,\mathrm{d}x\,\mathrm{d}t\leqslant k\mu\int_{\mathbb{R}^{m}}\Bigl{[}\tfrac{1}{2}v_{0}^{2}+\sum_{\alpha=1}^{m}\Phi_{\alpha}(Z_{\alpha 0})\Bigr{]}\,\mathrm{d}x.
\label{align:ch06-6-6-18}
\end{align}
正如\cref{theorem:ch06-6-6-1}证明中所释，\eqref{align:ch06-6-6-1}的弱解由经典解序列的 L$^{1}$ 极限构造，故不等式\eqref{align:ch06-6-6-18}对初值属于 $L^{1}(\mathbb{R}^{m})\cap L^{\infty}(\mathbb{R}^{m})$ 的弱解也成立。

\begin{theorem}\label{theorem:ch06-6-6-2}
设 $(v^{\mu},Z_{1}^{\mu},\cdots,Z_{m}^{\mu})$ 表示\eqref{align:ch06-6-6-1}、\eqref{align:ch06-6-6-6}以参数 $\mu>0$、初值 $(v_{0},F_{1}(v_{0}),\cdots,F_{m}(v_{0}))$（$v_{0}\in L^{1}(\mathbb{R}^{m})\cap L^{\infty}(\mathbb{R}^{m})$）的解族。则存在 $\mathbb{R}^{m}\times[0,\infty)$ 上的有界可测函数 $v$，使当 $\mu\downarrow 0$ 时，
\begin{align}
v^{\mu}(x,t)\longrightarrow v(x,t),\qquad Z_{\alpha}^{\mu}(x,t)\longrightarrow F_{\alpha}(v(x,t)),\ \alpha=1,\cdots,m,
\label{align:ch06-6-6-19}
\end{align}
在 $\mathbb{R}^{m}\times[0,\infty)$ 上几乎处处成立。函数
\begin{align}
u(x,t)=v(x,t)-\sum_{\alpha=1}^{m}F_{\alpha}(v(x,t))
\label{align:ch06-6-6-20}
\end{align}
是守恒律\eqref{align:ch06-6-1-1}（流通量 $G_{\alpha}$ 由\eqref{align:ch06-6-6-5}确定）以初值
\begin{align}
u_{0}(x)=v_{0}(x)-\sum_{\alpha=1}^{m}F_{\alpha}(v_{0}(x)),\qquad x\in\mathbb{R}^{m}
\label{align:ch06-6-6-21}
\end{align}
为初值的可容许弱解。
\end{theorem}

\begin{proof}
对 $(x,t)\in\mathbb{R}^{m}\times[0,\infty)$，令
\begin{align}
u^{\mu}(x,t)=v^{\mu}(x,t)-\sum_{\alpha=1}^{m}Z_{\alpha}^{\mu}(x,t),
\label{align:ch06-6-6-22}
\end{align}
\begin{align}
G_{\alpha}^{\mu}(x,t)=c_{\alpha}[v^{\mu}(x,t)+Z_{\alpha}^{\mu}(x,t)].
\label{align:ch06-6-6-23}
\end{align}
由\eqref{align:ch06-6-6-4}，
\begin{align}
\partial_{t}u^{\mu}(x,t)+\operatorname{div}G^{\mu}(x,t)=0.
\label{align:ch06-6-6-24}
\end{align}

先证存在 $\mathbb{R}^{m}\times[0,\infty)$ 上的有界可测函数 $u$ 及某序列 $\{\mu_{n}\}$（$\mu_{n}\downarrow 0$，$n\to\infty$），使对一切 $t\in[0,\infty)$，
\begin{align}
u^{\mu_{n}}(\cdot,t)\longrightarrow u(\cdot,t),\qquad n\to\infty,
\label{align:ch06-6-6-25}
\end{align}
在 $L^{\infty}(\mathbb{R}^{m})$ 弱$^{*}$拓扑下。为此，固定 $C_{0}^{\infty}(\mathbb{R}^{m})$ 中任一试验函数 $\chi$，定义函数族
\begin{align}
w^{\mu}(t)=\int_{\mathbb{R}^{m}}\chi(x)u^{\mu}(x,t)\,\mathrm{d}x,\qquad t\in[0,\infty).
\label{align:ch06-6-6-26}
\end{align}
由\eqref{align:ch06-6-6-24}，$w^{\mu}$ 连续可微，其导数
\begin{align}
\frac{\mathrm{d}}{\mathrm{d}t}w^{\mu}(t)=\sum_{\alpha=1}^{m}\int_{\mathbb{R}^{m}}\partial_{\alpha}\chi(x)\,G_{\alpha}^{\mu}(x,t)
\label{align:ch06-6-6-27}
\end{align}
关于 $\mu>0$ 一致有界。于是由 Arzelà 定理，存在序列 $\{\mu_{n}\}$（$\mu_{n}\downarrow 0$，$n\to\infty$），使 $\{w^{\mu_{n}}\}$ 对一切 $t\in[0,\infty)$ 收敛。通过 Cantor 对角线程序，可构造 $\{\mu_{n}\}$ 的子列（仍记为 $\{\mu_{n}\}$），使序列
\begin{align}
\Bigl{\{}\int_{\mathbb{R}^{m}}\chi(x)u^{\mu_{n}}(x,t)\,\mathrm{d}x\Bigr{\}}
\label{align:ch06-6-6-28}
\end{align}
对一切 $t\in[0,\infty)$ 以及任一给定可数试验函数族中每个 $\chi$ 收敛。因而\eqref{align:ch06-6-6-28}对任一 $\chi\in L^{1}(\mathbb{R}^{m})$ 收敛。于是对每个 $t\in[0,\infty)$，存在 $\mathbb{R}^{m}$ 上有界可测函数（记为 $u(\cdot,t)$），使\eqref{align:ch06-6-6-25}在 $L^{\infty}(\mathbb{R}^{m})$ 弱$^{*}$拓扑下成立。再注意到，由 L$^{1}$ 压缩估计\eqref{align:ch06-6-6-11}，对任一固定 $t\in[0,\infty)$，函数族 $(v^{\mu}(\cdot,t),Z_{1}^{\mu}(\cdot,t),\cdots,Z_{m}^{\mu}(\cdot,t))$ 在平均意义下等度连续。因而\eqref{align:ch06-6-6-25}中的收敛升级为在 $L^{1}(\mathbb{R}^{m})$ 中强收敛。特别地，
\begin{align}
u^{\mu_{n}}(x,t)\longrightarrow u(x,t),\qquad n\to\infty,
\label{align:ch06-6-6-29}
\end{align}
在 $\mathbb{R}^{m}\times[0,\infty)$ 上几乎处处成立。

对解 $(v^{\mu_{n}},Z_{1}^{\mu_{n}},\cdots,Z_{m}^{\mu_{n}})$ 应用\eqref{align:ch06-6-6-18}，必要时再取子列（仍记为 $\{\mu_{n}\}$），得
\begin{align}
F_{\alpha}(v^{\mu_{n}}(x,t))-Z_{\alpha}^{\mu_{n}}(x,t)\to 0,\qquad n\to\infty,\ \alpha=1,\cdots,m,
\label{align:ch06-6-6-30}
\end{align}
在 $\mathbb{R}^{m}\times[0,\infty)$ 上几乎处处成立。

结合\eqref{align:ch06-6-6-22}、\eqref{align:ch06-6-6-29}与\eqref{align:ch06-6-6-30}，得
\begin{align}
v^{\mu_{n}}(x,t)-\sum_{\alpha=1}^{m}F_{\alpha}(v^{\mu_{n}}(x,t))\to u(x,t),\qquad n\to\infty,
\label{align:ch06-6-6-31}
\end{align}
在 $\mathbb{R}^{m}\times[0,\infty)$ 上几乎处处成立。由单调性假设\eqref{align:ch06-6-6-2}，\eqref{align:ch06-6-6-31}蕴含序列 $\{v^{\mu_{n}}\}$ 本身必收敛，设
\begin{align}
v^{\mu_{n}}(x,t)\to v(x,t),\qquad n\to\infty,
\label{align:ch06-6-6-32}
\end{align}
在 $\mathbb{R}^{m}\times[0,\infty)$ 上几乎处处成立，其中 $v$ 是与 $u$ 通过\eqref{align:ch06-6-6-20}相联系的函数。此外，\eqref{align:ch06-6-6-30}与\eqref{align:ch06-6-6-32}共同蕴含
\begin{align}
Z_{\alpha}^{\mu_{n}}(x,t)\to F_{\alpha}(v(x,t)),\qquad n\to\infty,\ \alpha=1,\cdots,m,
\label{align:ch06-6-6-33}
\end{align}
在 $\mathbb{R}^{m}\times[0,\infty)$ 上几乎处处成立。

由\eqref{align:ch06-6-6-22}、\eqref{align:ch06-6-6-23}、\eqref{align:ch06-6-6-24}、\eqref{align:ch06-6-6-32}与\eqref{align:ch06-6-6-33}，$u$ 是\eqref{align:ch06-6-1-1}以 $G_{\alpha}$ 由\eqref{align:ch06-6-6-5}确定的弱解。下面证明此解可容许。固定任一常数 $\bar{v}$，对解 $(v^{\mu_{n}},Z_{1}^{\mu_{n}},\cdots,Z_{m}^{\mu_{n}})$ 与 $(\bar{v},F_{1}(\bar{v}),\cdots,F_{m}(\bar{v}))$ 写出\eqref{align:ch06-6-6-14}。把此（分布）方程应用于 $\mathbb{R}^{m}\times[0,\infty)$ 上任一具紧支集的非负 Lipschitz 连续试验函数 $\psi$，令 $\varepsilon\downarrow 0$。由于\eqref{align:ch06-6-6-14}右端当 $\varepsilon\downarrow 0$ 时非正，此计算给出
\begin{align}
&\int_{0}^{\infty}\int_{\mathbb{R}^{m}}\partial_{t}\psi\Bigl{[}(v^{\mu_{n}}-\bar{v})^{+}+\sum_{\alpha=1}^{m}(F_{\alpha}(\bar{v})-Z_{\alpha}^{\mu_{n}})^{+}\Bigr{]}\,\mathrm{d}x\,\mathrm{d}t\notag\\
&\qquad+\int_{0}^{\infty}\int_{\mathbb{R}^{m}}\sum_{\alpha=1}^{m}c_{\alpha}\partial_{\alpha}\psi\Bigl{[}(v^{\mu_{n}}-\bar{v})^{+}-(F_{\alpha}(\bar{v})-Z_{\alpha}^{\mu_{n}})^{+}\Bigr{]}\,\mathrm{d}x\,\mathrm{d}t\notag\\
&\qquad+\int_{\mathbb{R}^{m}}\psi(x,0)\Bigl{[}(v_{0}-\bar{v})^{+}+\sum_{\alpha=1}^{m}(F_{\alpha}(\bar{v})-F_{\alpha}(v_{0}))^{+}\Bigr{]}\,\mathrm{d}x\geqslant 0.
\label{align:ch06-6-6-34}
\end{align}
令 $n\to\infty$，利用\eqref{align:ch06-6-6-32}与\eqref{align:ch06-6-6-33}，\eqref{align:ch06-6-6-34}给出
\begin{align}
&\int_{0}^{\infty}\int_{\mathbb{R}^{m}}\partial_{t}\psi\Bigl{[}(v-\bar{v})^{+}+\sum_{\alpha=1}^{m}(F_{\alpha}(\bar{v})-F_{\alpha}(v))^{+}\Bigr{]}\,\mathrm{d}x\,\mathrm{d}t\notag\\
&\qquad+\int_{0}^{\infty}\int_{\mathbb{R}^{m}}\sum_{\alpha=1}^{m}c_{\alpha}\partial_{\alpha}\psi\Bigl{[}(v-\bar{v})^{+}-(F_{\alpha}(\bar{v})-F_{\alpha}(v))^{+}\Bigr{]}\,\mathrm{d}x\,\mathrm{d}t\notag\\
&\qquad+\int_{\mathbb{R}^{m}}\psi(x,0)\Bigl{[}(v_{0}-\bar{v})^{+}+\sum_{\alpha=1}^{m}(F_{\alpha}(\bar{v})-F_{\alpha}(v_{0}))^{+}\Bigr{]}\,\mathrm{d}x\geqslant 0.
\label{align:ch06-6-6-35}
\end{align}
由\eqref{align:ch06-6-6-2}，$v-\bar{v}$ 与 $F_{\alpha}(\bar{v})-F_{\alpha}(v)$ 同号。此外，令 $\bar{u}=\bar{v}-\sum_{\alpha=1}^{m}F_{\alpha}(\bar{v})$，则 $v-\bar{v}$ 与 $u-\bar{u}$ 也同号。于是利用\eqref{align:ch06-6-6-20}、\eqref{align:ch06-6-6-21}与\eqref{align:ch06-6-6-5}，可把\eqref{align:ch06-6-6-35}改写为
\begin{align}
\int_{0}^{\infty}\int_{\mathbb{R}^{m}}\Bigl{[}\partial_{t}\psi\,\eta(u;\bar{u})+\sum_{\alpha=1}^{m}\partial_{\alpha}\psi\,Q_{\alpha}(u;\bar{u})\Bigr{]}\,\mathrm{d}x\,\mathrm{d}t+\int_{\mathbb{R}^{m}}\psi(x,0)\eta(u_{0};\bar{u})\,\mathrm{d}x\geqslant 0,
\label{align:ch06-6-6-36}
\end{align}
其中 $(\eta(u;\bar{u}),Q(u;\bar{u}))$ 是\eqref{align:ch06-6-2-5}所定义的熵--熵流对。正如\cref{section:ch06-可容许弱解稳定性}所注，以 $\bar{u}$ 任意的熵--熵流对族\eqref{align:ch06-6-2-5}是“完备”的，因而\eqref{align:ch06-6-6-36}蕴含对任一凸熵--熵流对 $(\eta,Q)$，\eqref{align:ch06-6-2-3}成立。这验证了 $u$ 是以\eqref{align:ch06-6-6-21}所给 $u_{0}$ 为初值的\eqref{align:ch06-6-1-1}的可容许弱解。由于 $u$ 唯一，\eqref{align:ch06-6-6-29}、\eqref{align:ch06-6-6-32}与\eqref{align:ch06-6-6-33}中的收敛不仅沿子列 $\{\mu_{n}\}$ 成立，而且沿整个族 $\{\mu\}$ 当 $\mu\downarrow 0$ 成立。证明完成。

\end{proof}

\cref{theorem:ch06-6-6-2}表明，从给定系统\eqref{align:ch06-6-6-1}出发，如何通过松弛构造由\eqref{align:ch06-6-6-1}诱导的某一标量守恒律的可容许解。当然，我们感兴趣的是反过程：确定适当的系统\eqref{align:ch06-6-6-1}，使其松弛形式为给定的标量守恒律\eqref{align:ch06-6-1-1}。当给定流通量 $G_{\alpha}(u)$ 时，可通过选取常数 $c_{\alpha}$ 使变换\eqref{align:ch06-6-6-5}隐式确定的函数 $F_{\alpha}(v)$ 满足假设\eqref{align:ch06-6-6-2}与\eqref{align:ch06-6-6-3}。归一化给定流通量使 $G_{\alpha}(0)=0$，$\alpha=1,\cdots,m$。由于解先验有界，可不失一般性地设 $G_{\alpha}^{\prime}(u)$ 在 $(-\infty,\infty)$ 上一致有界。由\eqref{align:ch06-6-6-5}，
\begin{align}
(m+1)v=u+\sum_{\alpha=1}^{m}\frac{1}{c_{\alpha}}G_{\alpha}(u).
\label{align:ch06-6-6-37}
\end{align}
因而第一个约束是把 $|c_{\alpha}|$ 固定得充分大，使
\begin{align}
(m+1)\frac{\mathrm{d}v}{\mathrm{d}u}=1+\sum_{\alpha=1}^{m}\frac{1}{c_{\alpha}}G_{\alpha}^{\prime}(u)\geqslant\frac{1}{2},
\label{align:ch06-6-6-38}
\end{align}
以保证映射 $v\mapsto u$ 具有光滑逆。再注意
\begin{align}
F_{\alpha}^{\prime}(v)=-1+\frac{1}{c_{\alpha}}G_{\alpha}^{\prime}(u)\frac{\mathrm{d}u}{\mathrm{d}v}=-1+\frac{m+1}{c_{\alpha}}\Bigl{[}1+\sum_{\beta=1}^{m}\frac{1}{c_{\beta}}G_{\beta}^{\prime}(u)\Bigr{]}^{-1}G_{\alpha}^{\prime}(u),
\label{align:ch06-6-6-39}
\end{align}
于是通过把 $|c_{\alpha}|$ 选得充分大，可同时满足假设\eqref{align:ch06-6-6-2}与\eqref{align:ch06-6-6-3}。要求对流特征速度 $c_{\alpha}$ 相对于松弛守恒律的特征速度 $G_{\alpha}^{\prime}$ 足够大的对 $c_{\alpha}$ 的限制，称为\emph{次特征条件}。

\end{document}
